% GENERATED FILE. Edit canonical lessons, then run npm run generate:books.
% canonical-source-hash: fnv1a64:37b8c6054a7f5b82
% canonical-lessons: ZH-C154-hai, ZH-W154-hai-still, ZH-C154-haishi, ZH-C154-huozhe, ZH-C154-bie, ZH-C154-practice

\chapter{This One or That One}
\label{ch:zh-this-one-or-that-one}
\hlchaptermodality{154}

\begin{chapteropening}
\textbf{By the end of this chapter:} \emph{I can ask 'this or that?', say 'not yet' and 'don't'.}

Still and not yet with \zh{还}, asking which of two with \zh{还是}, offering either with huòzhě, and don't with bié, built on one new written character: \zh{还}.
\end{chapteropening}

\section[\texorpdfstring{\zh{还} (hái)}{还 (hái)}]{\zh{还} (hái) — still; yet}

\label{lesson:ZH-C154-hai}



\begin{quote}
A new chapter, and it opens on a word for what has not changed yet.

\begin{itemize}
\raggedright
  \item \emph{Recall:} say \zh{我比他大} — \emph{wǒ bǐ tā dà} — \textbf{R1}, the chapter you have finished
  \item \emph{Recall:} say \zh{那么} — \emph{nàme} — \textbf{R1}
  \item \emph{Recall:} say \emph{hái xíng}, not bad — \textbf{R4}, recognition at distance
  \item \emph{Recall:} say \emph{yǐjīng}, already — \textbf{R2}, a few chapters back
\end{itemize}
\end{quote}

\subsection*{You'll want to know: \zh{还}}
\textbf{\zh{还}} — \emph{hái} — \textquotedblleft{}still; yet\textquotedblright{}.

\begin{quote}
\textbf{\zh{他还在家}} \emph{tā hái zài jiā} — He's still at home.
\end{quote}

\begin{quote}
\textbf{\zh{我还没吃饭}} \emph{wǒ hái méi chī fàn} — I haven't eaten yet.
\end{quote}

The next lesson puts the character in your hand. It is the hái of \emph{hái xíng}.

\begin{cousinweb}
\textbf{\zh{还}} (\emph{hái}) is the hái of \emph{hái xíng}, inherited from Middle Chinese \emph{hwaen} (Guangyun, in Baxter's transcription), glossed 'to return'. The regular sound change gives huán, the \emph{huán} 'to give back' you know, not hái; how hái split off from it, and how it came to mean 'still', are not settled.
\end{cousinweb}

\begin{grammarlens}[title={What you've built}]
\zh{还} goes before the verb: still. \zh{还没} before a verb is not yet: \zh{我还没吃饭}.

Set it against \emph{yǐjīng}, already: \emph{wǒ yǐjīng chī le}, I've eaten already; \zh{我还没吃饭}, not yet.
\end{grammarlens}

\subsection*{Your turn}
\begin{itemize}
\raggedright
  \item \emph{Say it:} \emph{hái}
  \item \emph{Say it:} \emph{hái}, once more
  \item \emph{Say it:} \zh{他还在家} — \emph{tā hái zài jiā} — He's still at home.
  \item \emph{Say it:} \zh{我还没吃饭} — \emph{wǒ hái méi chī fàn} — I haven't eaten yet.
  \item \emph{Recall:} say \emph{yǐjīng}, then \zh{还没}
\end{itemize}

\begin{tcolorbox}[breakable,colback=teal!4,colframe=teal!35!black,title={Before you move on}]
What does \zh{还} mean? (\textquotedblleft{}Still; yet\textquotedblright{}.) How do you say not yet? (\textbf{\zh{还没}, before the verb}.)
\end{tcolorbox}
\section[\texorpdfstring{\zh{还} (hái)}{还 (hái)}]{\zh{还} — \zh{不}, then the walking piece last}

\label{lesson:ZH-W154-hai-still}



\begin{quote}
The word first, then the character inside it.

\begin{itemize}
\raggedright
  \item \emph{Recall:} say \emph{hái} — \textbf{R1}, one lesson back, before it decays
\end{itemize}

\begin{quote}
\zh{不} \zh{这} \zh{过}
\end{quote}

\emph{Read it:} each one aloud, with nothing said first — \emph{bù}, \emph{zhè}, \emph{guò}
\end{quote}

\subsection*{Script — \zh{不}, then the walking piece last}
\hlblockfigure{\detokenize{figures/ZH-W154-hai-still-filmstrip.pdf}}{How \zh{还} is written, stroke by stroke}

\begin{quote}
\zh{还}
\end{quote}

\textbf{Two pieces you own, in the order of \zh{这} and \zh{过}.}

\textbf{\zh{不}, four strokes}, on the upper right: the bar, the left-falling stroke, the vertical, the dot.

\textbf{\zh{辶}, three strokes}, last: the dot, the stroke that turns down, turns again and falls to the left, then the long base under \zh{不}.

4 + 3 = seven.

\subsection*{Your turn}
\begin{itemize}
\raggedright
  \item \emph{Write it:} \zh{还}, all seven, with the model in front of you
  \item \emph{Cover:} write it again from memory, and count the pieces, not the strokes
\end{itemize}

\begin{tcolorbox}[breakable,colback=teal!4,colframe=teal!35!black,title={Before you move on}]
How many strokes in \zh{还}? (\textbf{Seven}.) Which piece goes in first? (\textbf{\zh{不}}.) Which other characters put the walking piece last? (\textbf{\zh{这}, \zh{进} and \zh{过}}.)

Source: \href{https://raw.githubusercontent.com/chanind/hanzi-writer-data/68d10a4b21150cae5e1ebbd223eed289cf32d90c/data/\%E8\%BF\%98.json}{Hanzi Writer Data \zh{还}, pinned stroke paths and medians}.
\end{tcolorbox}
\section[\texorpdfstring{\zh{还是} (háishi)}{还是 (háishi)}]{\zh{还是} (háishi) — or, in a question}

\label{lesson:ZH-C154-haishi}



\begin{quote}
Before the new one: read the character you have written, then a sentence with \zh{要}.

\begin{itemize}
\raggedright
  \item \emph{Read it:} \zh{还没}, aloud — \emph{hái méi}, not yet
  \item \emph{Recall:} say \zh{我要水} — \emph{wǒ yào shuǐ}, I want some water — \textbf{R3}, durable
\end{itemize}
\end{quote}

\subsection*{You'll want to know: \zh{还是}}
\textbf{\zh{还是}} — \emph{háishi} — \textquotedblleft{}or\textquotedblright{}, in a question.

\begin{quote}
\textbf{\zh{你要水还是要饭}} \emph{nǐ yào shuǐ háishi yào fàn} — Do you want water or food?
\end{quote}

\begin{quote}
\textbf{\zh{你去还是他去}} \emph{nǐ qù háishi tā qù} — Are you going, or is he?
\end{quote}

\begin{cousinweb}
\textbf{\zh{还是}} (\emph{háishi}) is built from \zh{还} 'still', the character you have written, and \zh{是} 'to be'. Both are inherited from Middle Chinese, \emph{hwaen} and \emph{dzyeX} (Guangyun, in Baxter's transcription); 'or' is the pair's own sense, which neither gloss gives.
\end{cousinweb}

\begin{grammarlens}[title={What you've built}]
\zh{还是} goes between the two choices, and the answer picks one: \zh{我要水}.

Each choice can keep its own verb: \zh{要水还是要饭}. The question needs no \zh{吗}: \zh{还是} already makes it one.
\end{grammarlens}

\subsection*{Your turn}
\begin{itemize}
\raggedright
  \item \emph{Say it:} \emph{háishi}
  \item \emph{Say it:} \emph{háishi}, once more
  \item \emph{Say it:} \zh{你要水还是要饭} — \emph{nǐ yào shuǐ háishi yào fàn} — Do you want water or food?
  \item \emph{Say it:} \zh{我要水} — \emph{wǒ yào shuǐ} — Water, please.
  \item \emph{Recall:} say \zh{还没}, then \zh{还是}, and hear the same hái
\end{itemize}

\begin{tcolorbox}[breakable,colback=teal!4,colframe=teal!35!black,title={Before you move on}]
What does \zh{还是} do? (\textbf{Asks which of two}.) Does the question need \zh{吗}? (\textbf{No}.)
\end{tcolorbox}
\section[huòzhě]{huòzhě — or, in a statement}

\label{lesson:ZH-C154-huozhe}



\begin{quote}
Before the new one: last lesson's sentence, then the or that asks.

\begin{itemize}
\raggedright
  \item \emph{Recall:} say \zh{你要水还是要饭} — \emph{nǐ yào shuǐ háishi yào fàn} — \textbf{R1}, one lesson back
  \item \emph{Recall:} say \zh{我还没吃饭} — \emph{wǒ hái méi chī fàn} — \textbf{R1}, this chapter
\end{itemize}
\end{quote}

\subsection*{You'll want to know: huòzhě}
\emph{huòzhě} — \textquotedblleft{}or\textquotedblright{}, in a statement.

\begin{quote}
\emph{míngtiān huòzhě hòutiān dōu kěyǐ} — Tomorrow or the day after: either is fine.
\end{quote}

Its characters come later; for now it is a word for the ear.

\begin{cousinweb}
\emph{huòzhě} is built from huò and zhě. Both are inherited from Middle Chinese, \emph{hwok} and \emph{tsyaeX} (Guangyun, in Baxter's transcription), glossed 'not fixed; in doubt' and 'a particle'. 'Or' for the pair is assumed here, not quoted.
\end{cousinweb}

\begin{grammarlens}[title={What you've built}]
\zh{还是} asks; \emph{huòzhě} states. \textbf{\zh{你去还是他去}} \emph{nǐ qù háishi tā qù} wants an answer; \emph{nǐ qù huòzhě tā qù dōu kěyǐ} says either will do.

\zh{都} after the two choices says both are fine.
\end{grammarlens}

\subsection*{Your turn}
\begin{itemize}
\raggedright
  \item \emph{Say it:} \emph{huòzhě}
  \item \emph{Say it:} \emph{huòzhě}, once more
  \item \emph{Say it:} \emph{míngtiān huòzhě hòutiān dōu kěyǐ} — Tomorrow or the day after: either is fine.
  \item \emph{Say it:} \zh{你去还是他去} — \emph{nǐ qù háishi tā qù} — Are you going, or is he?
  \item \emph{Recall:} say \zh{还是}, then \emph{huòzhě}: one asks, one states
\end{itemize}

\begin{tcolorbox}[breakable,colback=teal!4,colframe=teal!35!black,title={Before you move on}]
Which or asks, \zh{还是} or \emph{huòzhě}? (\textbf{\zh{还是}}.) Which one states? (\emph{\textbf{huòzhě}}.)
\end{tcolorbox}
\section[bié]{bié — don't}

\label{lesson:ZH-C154-bie}



\begin{quote}
Before the new one: last lesson's word, then the don't you already know.

\begin{itemize}
\raggedright
  \item \emph{Recall:} say \emph{huòzhě}, or — \textbf{R1}, one lesson back
  \item \emph{Recall:} say \zh{不要说} — \emph{bú yào shuō}, don't say it — \textbf{R3}, durable
  \item \emph{Recall:} say \emph{biérén}, other people — \textbf{R2}
\end{itemize}
\end{quote}

\subsection*{You'll want to know: bié}
\emph{bié} — \textquotedblleft{}don't\textquotedblright{}.

\begin{quote}
\emph{bié qù} — Don't go.
\end{quote}

\begin{quote}
\emph{bié shuō le} — Don't say any more.
\end{quote}

Its character comes later; for now it is a word for the ear.

\begin{cousinweb}
\emph{bié} is inherited from Middle Chinese \emph{bjet} (Guangyun, in Baxter's transcription), glossed 'different; to part'. It is the bié of \emph{biérén} and \emph{biéde}. How it came to mean 'don't' is not settled.
\end{cousinweb}

\begin{grammarlens}[title={What you've built}]
\emph{bié} and \zh{不要} both stop someone: \emph{bié qù} and \zh{不要去} say the same thing, and \emph{bié} is shorter.

With \zh{了} at the end it means stop now, enough: \emph{bié shuō le}.
\end{grammarlens}

\subsection*{Your turn}
\begin{itemize}
\raggedright
  \item \emph{Say it:} \emph{bié}
  \item \emph{Say it:} \emph{bié}, once more
  \item \emph{Say it:} \emph{bié qù} — Don't go.
  \item \emph{Say it:} \zh{不要去} — \emph{bú yào qù} — Don't go.
  \item \emph{Say it:} \emph{bié shuō le} — Don't say any more.
  \item \emph{Recall:} say \zh{不要说}, then \emph{bié shuō}
\end{itemize}

\begin{tcolorbox}[breakable,colback=teal!4,colframe=teal!35!black,title={Before you move on}]
What does \emph{bié} mean? (\textquotedblleft{}Don't\textquotedblright{}.) Which longer words say the same? (\textbf{\zh{不要}}.)
\end{tcolorbox}
\section[Practice]{\zh{你要水还是要饭}}

\label{lesson:ZH-C154-practice}



\begin{quote}
The chapter closes here, with two people choosing between things.

\begin{itemize}
\raggedright
  \item \emph{Recall:} say \emph{bié}, don't — \textbf{R1}, one lesson back
  \item \emph{Recall:} say \emph{huòzhě}, or — \textbf{R1}
  \item \emph{Recall:} say \zh{还没} — \emph{hái méi}, not yet — \textbf{R2}
\end{itemize}
\end{quote}

\subsection*{The exchange}
\begin{quote}
\textbf{A:} \zh{你要水还是要饭}
\end{quote}

\begin{quote}
\emph{nǐ yào shuǐ háishi yào fàn} — Do you want water or food?
\end{quote}

\begin{quote}
\textbf{B:} \zh{我要水}
\end{quote}

\begin{quote}
\emph{wǒ yào shuǐ} — Water, please.
\end{quote}

\begin{quote}
\textbf{A:} \zh{你吃饭了吗}
\end{quote}

\begin{quote}
\emph{nǐ chī fàn le ma} — Have you eaten?
\end{quote}

\begin{quote}
\textbf{B:} \zh{还没}
\end{quote}

\begin{quote}
\emph{hái méi} — Not yet.
\end{quote}

\begin{quote}
\textbf{A:} \zh{我们去吃饭还是在家吃}
\end{quote}

\begin{quote}
\emph{wǒmen qù chī fàn háishi zài jiā chī} — Shall we go out to eat, or eat at home?
\end{quote}

\begin{quote}
\textbf{B:} \zh{在家吃} \zh{妈妈还在做饭}
\end{quote}

\begin{quote}
\emph{zài jiā chī, māma hái zài zuò fàn} — At home. Mum's still cooking.
\end{quote}

\begin{quote}
\textbf{A:} \zh{我们明天去还是后天去}
\end{quote}

\begin{quote}
\emph{wǒmen míngtiān qù háishi hòutiān qù} — Are we going tomorrow or the day after?
\end{quote}

\begin{quote}
\textbf{B:} \emph{míngtiān huòzhě hòutiān dōu kěyǐ} — Tomorrow or the day after: either is fine.
\end{quote}

\begin{quote}
\textbf{A:} \emph{hǎo, bié wàng le} — Fine. Don't forget.
\end{quote}

\zh{还是} asks, and the answer picks one; \emph{huòzhě} refuses to pick. \zh{还没} is a whole answer: not yet. \emph{bié wàng le}: don't forget.

\begin{tcolorbox}[breakable,skin=enhanced,colback=orange!6,colframe=orange!45!black,boxrule=0.5pt,arc=1mm,left=6pt,right=6pt,top=4pt,bottom=4pt,fonttitle=\bfseries,title={Read this}]
\emph{Read it:} the exchange through once without stopping, then answer the question below
\end{tcolorbox}

\begin{grammarlens}[title={What you've built — or, still, not yet, don't}]
Four small words that steer a conversation:

\begin{itemize}
\raggedright
  \item \textbf{\zh{你要水还是要饭}} \emph{nǐ yào shuǐ háishi yào fàn} — water or food?
  \item \emph{míngtiān huòzhě hòutiān dōu kěyǐ} — tomorrow or the day after, either is fine.
  \item \textbf{\zh{我还没吃饭}} \emph{wǒ hái méi chī fàn} — I haven't eaten yet.
  \item \emph{bié qù} — don't go.
\end{itemize}

\textbf{Listen.} \emph{Nǐ hē chá háishi hē shuǐ?}

\textbf{Speak.} Offer a friend two choices, then tell them not to forget.

\emph{Write it:} \zh{还} — \zh{不} first, the walking piece last
\end{grammarlens}

\subsection*{Your turn}
\begin{itemize}
\raggedright
  \item \emph{Say it:} \zh{你要水还是要饭} — \emph{nǐ yào shuǐ háishi yào fàn} — Do you want water or food?
  \item \emph{Say it:} \zh{还没} — \emph{hái méi} — Not yet.
  \item \emph{Say it:} the whole exchange, both parts
  \item \emph{Write it:} \zh{你要水还是要饭}
\end{itemize}

\begin{tcolorbox}[breakable,colback=teal!4,colframe=teal!35!black,title={Before you move on}]
Ask water or food. (\textbf{\zh{你要水还是要饭}}.) Say not yet. (\textbf{\zh{还没}}.) Say either is fine. (\emph{míngtiān huòzhě hòutiān dōu kěyǐ}.) Say don't go. (\emph{bié qù}.)
\end{tcolorbox}
